% source: manti.01.ma, 1ère SM, devoir (Drive 0B6pfSJcEVng0VGIzSDBQeFF3NEU). Ex.1-2 trigonométrie / géométrie (hors thème)
\ex{Ensembles}{VGIz}{Ex. 3}{10}
On considère l'ensemble $A=\left\{\dfrac{x+y}{1+xy}\ /\ -1<x<1\ \text{et}\ -1<y<1\right\}$.
\begin{enumerate}
\item Montrer que $\left(\forall(x,y)\in\,]-1,1[^2\right)\ 1+xy\neq0$.
\item Montrer que $A\subset\,]-1,1[$.
\item Soit $a$ un réel de l'intervalle $]-1,1[$ ; on pose $x=\dfrac{a}{1+\sqrt{1-a^2}}$.
\begin{enumerate}
\item Montrer que $-1<x<1$.
\item Calculer $\dfrac{2x}{1+x^2}$. Que peut-on en déduire ?
\end{enumerate}
\end{enumerate}
\fin
\ex{Suites}{VGIz}{Ex. 4}{12}
\textbf{Partie 1.} On considère l'équation $(E)\ :\ x^3+6x-2=0$.
\begin{enumerate}
\item Tracer dans un même repère les courbes des fonctions $g(x)=x^3$ et $h(x)=2-6x$.
\item En déduire graphiquement que l'équation $(E)$ admet une unique solution $\alpha$.
\item \begin{enumerate}
\item Déterminer le signe de la fonction $F(x)=x^3+6x-2$.
\item En déduire que $\alpha$ appartient à l'intervalle $]0,1[$.
\end{enumerate}
\end{enumerate}
\textbf{Partie 2.} Soit $f$ la fonction définie par $f(x)=\dfrac16\left(2-x^3\right)$ et soit $I=[0,1]$.
\begin{enumerate}
\item Montrer que $f$ est décroissante sur $I$ et que $f(I)\subset I$.
\item Soit $(U_n)_n$ la suite définie par $U_0=\dfrac12$ et $U_{n+1}=f(U_n)$.
\begin{enumerate}
\item Montrer que $(\forall n\in\mathbb{N})\ 0<U_n<1$.
\item Vérifier que $f(\alpha)=\alpha$, puis montrer que $(\forall n\in\mathbb{N})\ |U_{n+1}-\alpha|\leq\dfrac12|U_n-\alpha|$.
\item En déduire que $(\forall n\in\mathbb{N})\ |U_n-\alpha|\leq\left(\dfrac12\right)^n$.
\end{enumerate}
\end{enumerate}
\fin
